% ============================================================================
\chapter{Preparing MNIST for Bitwise Processing}
\label{ch:mnist-prep}
% ============================================================================

Before any of our three paradigms can process MNIST~\cite{lecun1998mnist} images, the raw
pixel data must be converted to a Boolean representation. This chapter
covers the conversion methods and their trade-offs.


% ----------------------------------------------------------------------------
\section{The Raw Data}
\label{sec:raw-data}
% ----------------------------------------------------------------------------

Each MNIST image is a $28 \times 28$ grid of pixels. Each pixel is an
unsigned 8-bit integer in the range $[0, 255]$, where 0~is black
(background) and 255~is white (ink).

In practice, the images are stored as a flat array of 784 bytes:
\begin{equation}
  \bvec{x}_{\text{raw}} = (p_1, p_2, \ldots, p_{784}),
  \quad p_i \in \{0, 1, \ldots, 255\}
  \label{eq:raw-mnist}
\end{equation}

\begin{engineernote}
MNIST pixels are heavily bimodal: most pixels are either very dark
(background, $p \approx 0$) or quite bright (ink, $p > 200$). The
middle values ($50$--$150$) are relatively rare, occurring mainly at
the edges of strokes. This bimodal distribution is why simple
thresholding works well.
\end{engineernote}


% ----------------------------------------------------------------------------
\section{Simple Thresholding}
\label{sec:thresholding}
% ----------------------------------------------------------------------------

The simplest conversion to Boolean is \concept{thresholding}: compare
each pixel to a fixed threshold $\theta$ and output 1 if the pixel is
bright enough, 0 otherwise.

\begin{equation}
  x_k = \begin{cases}
    1 & \text{if } p_k > \theta \\
    0 & \text{if } p_k \leq \theta
  \end{cases}
  \label{eq:threshold}
\end{equation}

A common choice is $\theta = 127$ (the midpoint of the $[0, 255]$
range).

\begin{figure}[htbp]
\centering
\begin{tikzpicture}[
  px/.style={draw, minimum size=0.35cm, inner sep=0pt, font=\tiny\ttfamily}
]
  % Original
  \node[font=\small\sffamily\bfseries] at (2, 3.5) {Grayscale};
  \foreach \x/\y/\v in {
    0/2/0, 1/2/30, 2/2/180, 3/2/255,
    0/1/0, 1/1/200, 2/1/255, 3/1/150,
    0/0/0, 1/0/50, 2/0/220, 3/0/0
  }{
    \pgfmathsetmacro{\intensity}{\v/255*100}
    \node[px, fill=bitblue!\intensity] at (\x*0.5, \y*0.5) {\v};
  }

  % Arrow
  \draw[-{Stealth}, thick] (2.5, 0.75) -- (4, 0.75)
    node[midway, above, font=\scriptsize\sffamily] {$\theta = 127$};

  % Binary
  \node[font=\small\sffamily\bfseries] at (6, 3.5) {Boolean};
  \foreach \x/\y/\v in {
    0/2/0, 1/2/0, 2/2/1, 3/2/1,
    0/1/0, 1/1/1, 2/1/1, 3/1/1,
    0/0/0, 1/0/0, 2/0/1, 3/0/0
  }{
    \ifnum\v=1
      \def\fillcol{bitblue!40}
    \else
      \def\fillcol{white}
    \fi
    \node[px, fill=\fillcol] at (\x*0.5 + 4.5, \y*0.5) {\v};
  }
\end{tikzpicture}
\caption{Thresholding a $4 \times 3$ pixel patch at $\theta = 127$.
  Pixels above the threshold become~1 (foreground); others become~0
  (background).}
\label{fig:thresholding}
\end{figure}

After thresholding, the image is a bit string:
\begin{equation}
  \bvec{x} = (x_1, x_2, \ldots, x_{784}) \in \{0, 1\}^{784}
  \label{eq:bool-mnist}
\end{equation}


% ----------------------------------------------------------------------------
\section{Literal Expansion}
\label{sec:literal-expansion}
% ----------------------------------------------------------------------------

Tsetlin Machines~\cite{granmo2018tsetlin} need access to both a feature and its negation. For
every Boolean feature $x_k$, we also provide $\bar{x}_k = 1 - x_k$.

\begin{equation}
  \bvec{x}_{\text{expanded}} =
    (x_1, \bar{x}_1, x_2, \bar{x}_2, \ldots, x_{784}, \bar{x}_{784})
    \in \{0, 1\}^{1568}
  \label{eq:expanded}
\end{equation}

\begin{keyinsight}
Why include negations? Because a clause that says ``pixel~42 must be
dark'' needs the literal $\bar{x}_{42}$, which is 1 when the pixel is
dark (below threshold). Without negations, clauses could only express
``pixel must be bright'' conditions, halving their expressive power.
\end{keyinsight}

\begin{example}
For a $2 \times 2$ image with pixels $(200, 50, 180, 10)$ and
$\theta = 127$:
\begin{align*}
  \text{Thresholded:}\quad & (1, 0, 1, 0) \\
  \text{Expanded:}\quad & (1, 0, \; 0, 1, \; 1, 0, \; 0, 1)
\end{align*}
The expanded vector has 8 entries (twice the original). Notice that
$x_k$ and $\bar{x}_k$ are always complementary: exactly one of each
pair is~1.
\end{example}


% ----------------------------------------------------------------------------
\section{Thermometer Encoding}
\label{sec:thermometer}
% ----------------------------------------------------------------------------

Simple thresholding discards all information about pixel brightness
beyond ``above or below 127.'' A pixel value of 130 and a pixel value
of 255 both become~1, even though the latter is much brighter.

\concept{Thermometer encoding} uses multiple thresholds per pixel to
retain finer information:

\begin{definition}[Thermometer encoding]
Given $L$ threshold levels
$\theta_1 < \theta_2 < \cdots < \theta_L$, pixel $p_k$ is encoded as:
\begin{equation}
  t_k^{(j)} = \begin{cases}
    1 & \text{if } p_k > \theta_j \\
    0 & \text{otherwise}
  \end{cases}
  \quad \text{for } j = 1, \ldots, L
  \label{eq:thermometer}
\end{equation}
\end{definition}

\begin{example}
With thresholds $(50, 100, 150, 200)$, a pixel value of 175 encodes
as $(1, 1, 1, 0)$. A pixel value of 80 encodes as $(1, 0, 0, 0)$.
The encoding looks like a thermometer: 1s fill up from the bottom as
the value increases.
\end{example}

\begin{figure}[htbp]
\centering
\begin{tikzpicture}[
  bit/.style={draw, minimum width=0.6cm, minimum height=0.5cm,
    font=\scriptsize\ttfamily, inner sep=0pt}
]
  % Pixel value
  \node[font=\small\sffamily] at (-1.5, 0) {$p_k = 175$};

  % Thresholds
  \foreach \t/\i/\v in {50/0/1, 100/1/1, 150/2/1, 200/3/0} {
    \ifnum\v=1
      \def\fillcol{bitblue!30}
    \else
      \def\fillcol{white}
    \fi
    \node[bit, fill=\fillcol] (t\i) at (\i*0.8, 0) {\v};
    \node[font=\tiny\sffamily, bitgray, above] at (\i*0.8, 0.35) {$>\t$};
  }

  \draw[decorate, decoration={brace, amplitude=4pt, mirror}, thick]
    (-0.3, -0.5) -- (2.1, -0.5)
    node[midway, below=4pt, font=\scriptsize\sffamily] {4 bits per pixel};
\end{tikzpicture}
\caption{Thermometer encoding with 4 thresholds. The pixel value 175
  exceeds thresholds 50, 100, and 150 but not 200, yielding the code
  $(1, 1, 1, 0)$.}
\label{fig:thermometer}
\end{figure}

With $L$ thresholds per pixel, plus negations, the total input size is:
\begin{equation}
  \text{Input length} = 784 \times L \times 2
  \label{eq:thermo-size}
\end{equation}

For $L = 4$: $784 \times 4 \times 2 = 6{,}272$ literals. More
information is preserved, but the model must process a larger input.


% ----------------------------------------------------------------------------
\section{The Booleanisation Trade-Off}
\label{sec:bool-tradeoff}
% ----------------------------------------------------------------------------

\begin{table}[htbp]
\centering
\caption{Comparison of booleanisation strategies for MNIST.}
\label{tab:bool-comparison}
\begin{tabular}{@{}lccc@{}}
  \toprule
  \textbf{Method} & \textbf{Bits/pixel} & \textbf{Total literals}
    & \textbf{Info preserved} \\
  \midrule
  Simple threshold ($\theta = 127$) & 1 & 1{,}568 & Low \\
  Thermometer ($L = 4$) & 4 & 6{,}272 & Medium \\
  Thermometer ($L = 8$) & 8 & 12{,}544 & High \\
  Full 8-bit binary    & 8 & 12{,}544 & Full \\
  \bottomrule
\end{tabular}
\end{table}

\begin{engineernote}
For the proof-of-concept implementation described in this book, we
recommend starting with simple thresholding ($\theta = 127$). It
produces 1{,}568 literals, keeps the model small and fast, and achieves
accuracy within 1--2\% of the more sophisticated encodings. Once the
architecture is working, thermometer encoding can be added as an
optimisation.
\end{engineernote}


% ----------------------------------------------------------------------------
\section{Worked Example: One MNIST Digit Through Each Encoding}
\label{sec:encoding-walkthrough}
% ----------------------------------------------------------------------------

Let us trace a single pixel row from an MNIST ``3'' through the
encoding pipeline. Consider pixels at row 14 (the middle of the
image), positions 10--17, with raw values:

\begin{center}
\begin{tabular}{@{}l*{8}{c}@{}}
  \toprule
  \textbf{Position} & 10 & 11 & 12 & 13 & 14 & 15 & 16 & 17 \\
  \midrule
  Raw value & 0 & 42 & 190 & 255 & 255 & 180 & 30 & 0 \\
  \midrule
  Threshold ($\theta = 127$) & 0 & 0 & 1 & 1 & 1 & 1 & 0 & 0 \\
  \midrule
  Thermo ($>50$) & 0 & 0 & 1 & 1 & 1 & 1 & 0 & 0 \\
  Thermo ($>100$) & 0 & 0 & 1 & 1 & 1 & 1 & 0 & 0 \\
  Thermo ($>150$) & 0 & 0 & 1 & 1 & 1 & 1 & 0 & 0 \\
  Thermo ($>200$) & 0 & 0 & 0 & 1 & 1 & 0 & 0 & 0 \\
  \bottomrule
\end{tabular}
\end{center}

The simple threshold captures the stroke location (positions 12--15
are ink). The thermometer encoding additionally reveals that positions
13--14 are the \emph{darkest} part of the stroke (exceeding even the
200 threshold), while positions 12 and~15 are lighter edges.

After literal expansion, position~12 with simple thresholding
contributes two literals: $x_{12} = 1$ and $\bar{x}_{12} = 0$.
With 4-level thermometer encoding, it contributes 8 literals:
$(1, 0, 1, 0, 1, 0, 0, 1)$ for the four thresholds and their
negations.
